% Teoría de la partición espacial reproducible (modelo/clima.py), estilo cyberpunk.
% Compilar desde la raíz del proyecto:
%   python docs/generar_datos_latex.py      (regenera datos_generados.tex)
%   cd docs && pdflatex teoria_particion_espacial.tex && pdflatex teoria_particion_espacial.tex
% Versión para imprimir (fondo blanco): \usepackage[impresion]{estilo_cyberpunk}
\documentclass[11pt,a4paper]{article}
\usepackage[margin=2.4cm,headsep=0.8cm]{geometry}
\usepackage[utf8]{inputenc}
\usepackage[T1]{fontenc}
\usepackage{lmodern}
\usepackage[spanish,es-noshorthands,es-tabla,es-nolists]{babel}
\usepackage{amsmath,amssymb}
\usepackage{booktabs}
\usepackage{microtype}
\usepackage{pgfplots}
\pgfplotsset{compat=1.18}
\usepackage{estilo_cyberpunk}
\usepackage{hyperref}

\newcommand{\codigo}[1]{\cpcodigo{#1}}
\newcommand{\vect}[1]{\mathbf{#1}}
\newcommand{\trainset}{S_{\mathrm{train}}}
\newcommand{\valset}{S_{\mathrm{val}}}
\newcommand{\testset}{S_{\mathrm{test}}}

% Archivo generado por docs/generar_datos_latex.py; no editar a mano.

\newcommand{\TablaClimaSitios}{%
\begin{table}[htbp]\centering\scriptsize
\caption{Climatología NASA POWER 2001--2025 interpolada a los sitios de Ecuador y temporadas resultantes. $P$ y $ET_0$ son totales anuales (mm); $z_c$ es la elevación media de la celda MERRA-2 (m), no la del sitio.}\label{tab:clima-sitios}
\setlength{\tabcolsep}{2.5pt}
\begin{tabular}{@{}lllrrrlp{2.5cm}p{2.5cm}@{}}\toprule
Sitio & Provincia & Región & $P$ & $ET_0$ & $z_c$ & Régimen & Meses lluviosos & Meses secos ($<60$ mm) \\ \midrule
Humedal La Segua & Manabí & Costa & 947 & 1585 & 129 & unimodal & ene, feb, mar, abr & jun, jul, ago, sep, oct, nov, dic \\
Bosques andinos de Loja & Loja & Sierra & 662 & 1446 & 1818 & unimodal & ene, feb, mar, abr & may, jun, jul, ago, sep, oct, nov, dic \\
Pedernales & Manabí & Costa & 1166 & 1400 & 136 & unimodal & ene, feb, mar, abr, may & jun, jul, ago, sep, oct, nov \\
Glaciares del volcán Antisana & Napo & Amazonía & 1423 & 976 & 3121 & unimodal & ene, feb, mar, abr, may, nov, dic & ago \\
Arrozales de Daule & Guayas & Costa & 801 & 1972 & 100 & unimodal & ene, feb, mar, abr & may, jun, jul, ago, sep, oct, nov, dic \\
Guayaquil & Guayas & Costa & 751 & 1837 & 237 & unimodal & ene, feb, mar, abr & may, jun, jul, ago, sep, oct, nov, dic \\
\bottomrule\end{tabular}\end{table}}

\newcommand{\FiguraClimaSitios}{%
\begin{figure}[htbp]\centering
\begin{tikzpicture}
\begin{groupplot}[group style={group size=3 by 2, horizontal sep=1.1cm, vertical sep=1.4cm},
  width=0.36\textwidth, height=4.2cm, ymin=0, xmin=0.5, xmax=12.5, xtick={1,...,12},
  xticklabels={E,F,M,A,M,J,J,A,S,O,N,D}, tick label style={font=\scriptsize},
  title style={font=\small}, ylabel style={font=\scriptsize}, ymajorgrids]
\nextgroupplot[title={Humedal La Segua}, ylabel={mm/mes}]
\addplot[ybar, bar width=4.2pt, fill=lluviosa, draw=none] coordinates {(1,162.6) (2,215.1) (3,216.2) (4,153.2)};
\addplot[ybar, bar width=4.2pt, fill=transicion, draw=none] coordinates {(5,74.2)};
\addplot[ybar, bar width=4.2pt, fill=seca, draw=none] coordinates {(6,27.2) (7,10.0) (8,5.7) (9,7.8) (10,12.9) (11,14.9) (12,47.6)};
\addplot[thick, dashed, cplinea, mark=*, mark size=0.9pt] coordinates {(1,129.5) (2,108.5) (3,119.1) (4,117.8) (5,119.2) (6,117.1) (7,132.2) (8,148.6) (9,152.1) (10,154.9) (11,145.0) (12,141.1)};
\addplot[densely dotted, cpumbral, domain=0.5:12.5, samples=2] {60};
\nextgroupplot[title={Bosques andinos de Loja}]
\addplot[ybar, bar width=4.2pt, fill=lluviosa, draw=none] coordinates {(1,66.2) (2,99.5) (3,103.7) (4,70.8)};
\addplot[ybar, bar width=4.2pt, fill=seca, draw=none] coordinates {(5,44.8) (6,38.6) (7,31.2) (8,23.0) (9,22.9) (10,47.4) (11,53.5) (12,59.8)};
\addplot[thick, dashed, cplinea, mark=*, mark size=0.9pt] coordinates {(1,126.0) (2,110.6) (3,119.1) (4,112.5) (5,110.6) (6,100.7) (7,108.7) (8,124.4) (9,134.2) (10,139.2) (11,132.3) (12,128.1)};
\addplot[densely dotted, cpumbral, domain=0.5:12.5, samples=2] {60};
\nextgroupplot[title={Pedernales}]
\addplot[ybar, bar width=4.2pt, fill=lluviosa, draw=none] coordinates {(1,194.3) (2,222.8) (3,234.4) (4,185.0) (5,112.6)};
\addplot[ybar, bar width=4.2pt, fill=transicion, draw=none] coordinates {(12,69.3)};
\addplot[ybar, bar width=4.2pt, fill=seca, draw=none] coordinates {(6,50.4) (7,19.5) (8,11.3) (9,17.3) (10,24.3) (11,24.4)};
\addplot[thick, dashed, cplinea, mark=*, mark size=0.9pt] coordinates {(1,112.8) (2,96.6) (3,107.1) (4,104.5) (5,103.7) (6,100.9) (7,115.1) (8,131.4) (9,134.9) (10,138.1) (11,129.8) (12,124.7)};
\addplot[densely dotted, cpumbral, domain=0.5:12.5, samples=2] {60};
\nextgroupplot[title={Glaciares del volcán Antisana}, ylabel={mm/mes}]
\addplot[ybar, bar width=4.2pt, fill=lluviosa, draw=none] coordinates {(1,132.9) (2,137.8) (3,167.2) (4,178.3) (5,142.9) (11,126.6) (12,123.4)};
\addplot[ybar, bar width=4.2pt, fill=transicion, draw=none] coordinates {(6,98.8) (7,76.5) (9,66.2) (10,114.5)};
\addplot[ybar, bar width=4.2pt, fill=seca, draw=none] coordinates {(8,58.2)};
\addplot[thick, dashed, cplinea, mark=*, mark size=0.9pt] coordinates {(1,87.5) (2,79.6) (3,88.1) (4,81.1) (5,77.5) (6,67.3) (7,70.9) (8,78.1) (9,83.7) (10,90.4) (11,85.9) (12,85.4)};
\addplot[densely dotted, cpumbral, domain=0.5:12.5, samples=2] {60};
\nextgroupplot[title={Arrozales de Daule}]
\addplot[ybar, bar width=4.2pt, fill=lluviosa, draw=none] coordinates {(1,128.6) (2,216.6) (3,214.9) (4,131.4)};
\addplot[ybar, bar width=4.2pt, fill=seca, draw=none] coordinates {(5,42.3) (6,12.0) (7,4.3) (8,3.4) (9,3.0) (10,7.2) (11,9.5) (12,27.4)};
\addplot[thick, dashed, cplinea, mark=*, mark size=0.9pt] coordinates {(1,163.4) (2,132.4) (3,141.2) (4,142.8) (5,151.3) (6,152.3) (7,169.4) (8,186.4) (9,189.7) (10,190.3) (11,177.5) (12,175.3)};
\addplot[densely dotted, cpumbral, domain=0.5:12.5, samples=2] {60};
\nextgroupplot[title={Guayaquil}]
\addplot[ybar, bar width=4.2pt, fill=lluviosa, draw=none] coordinates {(1,114.6) (2,195.0) (3,200.2) (4,118.8)};
\addplot[ybar, bar width=4.2pt, fill=seca, draw=none] coordinates {(5,39.6) (6,13.7) (7,6.9) (8,5.7) (9,5.5) (10,11.5) (11,13.7) (12,26.1)};
\addplot[thick, dashed, cplinea, mark=*, mark size=0.9pt] coordinates {(1,153.0) (2,125.3) (3,133.2) (4,134.9) (5,142.7) (6,142.9) (7,157.4) (8,172.5) (9,175.5) (10,175.1) (11,163.0) (12,161.9)};
\addplot[densely dotted, cpumbral, domain=0.5:12.5, samples=2] {60};
\end{groupplot}
\end{tikzpicture}
\caption{Precipitación mensual climatológica (barras: cian lluviosa, gris transición, magenta seca), $ET_0$ de Hargreaves--Samani (línea discontinua) y umbral de mes seco de 60~mm (punteada) en los seis sitios del catálogo.}\label{fig:clima-sitios}
\end{figure}}

\newcommand{\FiguraForzamiento}{%
\begin{figure}[htbp]\centering
\begin{tikzpicture}
\begin{groupplot}[group style={group size=1 by 2, vertical sep=0.9cm}, width=0.92\textwidth,
  height=4.3cm, xmin=0, xmax=365, xtick={0,31,59,90,120,151,181,212,243,273,304,334},
  xticklabels={ene,feb,mar,abr,may,jun,jul,ago,sep,oct,nov,dic}, ymajorgrids,
  tick label style={font=\scriptsize}, ylabel style={font=\small}, legend style={font=\scriptsize, at={(0.99,0.97)}, anchor=north east}]
\nextgroupplot[ylabel={\%}, ymin=0, ymax=105, xticklabels={}]
\addplot[thick, color=lluviosa] coordinates {(0,33.8) (3,38.2) (6,42.8) (9,47.9) (12,53.2) (15,58.8) (18,64.6) (21,70.4) (24,76.3) (27,82.2) (30,88.3) (33,94.3) (36,100.0) (39,100.0) (42,100.0) (45,100.0) (48,100.0) (51,100.0) (54,100.0) (57,100.0) (60,100.0) (63,100.0) (66,100.0) (69,100.0) (72,100.0) (75,100.0) (78,100.0) (81,100.0) (84,100.0) (87,100.0) (90,100.0) (93,100.0) (96,100.0) (99,100.0) (102,100.0) (105,100.0) (108,100.0) (111,100.0) (114,100.0) (117,100.0) (120,99.8) (123,99.2) (126,98.1) (129,96.6) (132,94.7) (135,92.4) (138,90.0) (141,87.4) (144,84.7) (147,82.0) (150,79.1) (153,76.2) (156,73.2) (159,70.1) (162,67.0) (165,63.8) (168,60.6) (171,57.6) (174,54.6) (177,51.7) (180,48.9) (183,46.1) (186,43.4) (189,40.8) (192,38.3) (195,35.8) (198,33.5) (201,31.2) (204,29.2) (207,27.2) (210,25.3) (213,23.6) (216,21.9) (219,20.4) (222,19.0) (225,17.6) (228,16.3) (231,15.2) (234,14.1) (237,13.2) (240,12.4) (243,11.6) (246,11.0) (249,10.4) (252,9.8) (255,9.4) (258,8.9) (261,8.6) (264,8.3) (267,8.1) (270,7.9) (273,7.8) (276,7.7) (279,7.7) (282,7.7) (285,7.7) (288,7.7) (291,7.8) (294,7.9) (297,8.0) (300,8.1) (303,8.2) (306,8.3) (309,8.4) (312,8.6) (315,8.7) (318,8.8) (321,9.1) (324,9.5) (327,10.1) (330,10.8) (333,11.7) (336,12.7) (339,13.8) (342,15.0) (345,16.3) (348,17.7) (351,19.4) (354,21.6) (357,24.3) (360,27.5) (363,31.2)};
\addlegendentry{humedad del suelo $s$ (porcentaje del agua útil)}
\addplot[thick, color=seca] coordinates {(0,0.0) (3,0.0) (6,0.0) (9,0.0) (12,0.0) (15,0.0) (18,0.0) (21,0.0) (24,0.0) (27,0.0) (30,0.0) (33,0.0) (36,0.9) (39,12.9) (42,24.8) (45,36.5) (48,47.2) (51,56.5) (54,64.7) (57,71.8) (60,78.0) (63,83.3) (66,87.8) (69,91.7) (72,94.9) (75,97.5) (78,99.1) (81,99.9) (84,99.9) (87,99.2) (90,97.8) (93,95.8) (96,93.3) (99,90.3) (102,86.8) (105,83.0) (108,78.6) (111,73.7) (114,68.2) (117,62.3) (120,56.3) (123,50.9) (126,46.0) (129,41.5) (132,37.5) (135,33.9) (138,30.6) (141,27.6) (144,25.0) (147,22.5) (150,20.4) (153,18.4) (156,16.6) (159,15.0) (162,13.6) (165,12.2) (168,11.1) (171,10.0) (174,9.0) (177,8.2) (180,7.4) (183,6.7) (186,6.0) (189,5.4) (192,4.9) (195,4.4) (198,4.0) (201,3.6) (204,3.3) (207,2.9) (210,2.7) (213,2.4) (216,2.2) (219,2.0) (222,1.8) (225,1.6) (228,1.4) (231,1.3) (234,1.2) (237,1.1) (240,1.0) (243,0.9) (246,0.8) (249,0.7) (252,0.6) (255,0.6) (258,0.5) (261,0.5) (264,0.4) (267,0.4) (270,0.3) (273,0.3) (276,0.3) (279,0.3) (282,0.2) (285,0.2) (288,0.2) (291,0.2) (294,0.2) (297,0.1) (300,0.1) (303,0.1) (306,0.1) (309,0.1) (312,0.1) (315,0.1) (318,0.1) (321,0.1) (324,0.1) (327,0.1) (330,0.0) (333,0.0) (336,0.0) (339,0.0) (342,0.0) (345,0.0) (348,0.0) (351,0.0) (354,0.0) (357,0.0) (360,0.0) (363,0.0)};
\addlegendentry{almacenamiento del humedal $S$ (porcentaje del máximo anual)}
\nextgroupplot[ylabel={$\Delta\sigma^0_{suelo}$ (dB)}]
\addplot[thick, cpamarillo] coordinates {(0,-0.63) (3,-0.47) (6,-0.29) (9,-0.11) (12,0.09) (15,0.31) (18,0.52) (21,0.74) (24,0.96) (27,1.18) (30,1.41) (33,1.64) (36,1.85) (39,1.85) (42,1.85) (45,1.85) (48,1.85) (51,1.85) (54,1.85) (57,1.85) (60,1.85) (63,1.85) (66,1.85) (69,1.85) (72,1.85) (75,1.85) (78,1.85) (81,1.85) (84,1.85) (87,1.85) (90,1.85) (93,1.85) (96,1.85) (99,1.85) (102,1.85) (105,1.85) (108,1.85) (111,1.85) (114,1.85) (117,1.85) (120,1.84) (123,1.82) (126,1.78) (129,1.72) (132,1.65) (135,1.57) (138,1.47) (141,1.38) (144,1.28) (147,1.17) (150,1.07) (153,0.96) (156,0.84) (159,0.73) (162,0.61) (165,0.49) (168,0.37) (171,0.26) (174,0.15) (177,0.04) (180,-0.07) (183,-0.17) (186,-0.27) (189,-0.37) (192,-0.46) (195,-0.56) (198,-0.65) (201,-0.73) (204,-0.81) (207,-0.88) (210,-0.95) (213,-1.02) (216,-1.08) (219,-1.14) (222,-1.19) (225,-1.24) (228,-1.29) (231,-1.33) (234,-1.37) (237,-1.40) (240,-1.44) (243,-1.46) (246,-1.49) (249,-1.51) (252,-1.53) (255,-1.55) (258,-1.56) (261,-1.58) (264,-1.59) (267,-1.60) (270,-1.60) (273,-1.61) (276,-1.61) (279,-1.61) (282,-1.61) (285,-1.61) (288,-1.61) (291,-1.61) (294,-1.60) (297,-1.60) (300,-1.60) (303,-1.59) (306,-1.59) (309,-1.58) (312,-1.58) (315,-1.57) (318,-1.57) (321,-1.56) (324,-1.54) (327,-1.52) (330,-1.50) (333,-1.46) (336,-1.43) (339,-1.38) (342,-1.34) (345,-1.29) (348,-1.24) (351,-1.17) (354,-1.09) (357,-0.99) (360,-0.87) (363,-0.73)};
\end{groupplot}
\end{tikzpicture}
\caption{Ciclo anual del forzamiento estacional en el humedal La Segua (Manabí): humedad del suelo del balance hídrico, almacenamiento del humedal (embalse lineal, $\tau=30$ días) y cambio resultante de la retrodispersión del suelo, ecuación~\eqref{eq:dsigma}.}\label{fig:forzamiento}
\end{figure}}

\newcommand{\TablaHitos}{%
\begin{table}[htbp]\centering\small
\caption{Hitos astronómicos según la declinación de Spencer (día juliano $J$ y fecha de 2026) y, para cada sitio, pasos del Sol por el cenit ($\delta=\varphi$) y duración del día en los solsticios (h).}\label{tab:hitos}
\begin{tabular}{@{}lrr@{}}\toprule Hito & $J$ & Fecha \\ \midrule
Equinoccio de marzo & 80{,}16 & 21/03 \\
Solsticio de junio & 173{,}01 & 22/06 \\
Equinoccio de septiembre & 266{,}63 & 24/09 \\
Solsticio de diciembre & 356{,}29 & 22/12 \\
\bottomrule\end{tabular}\par\bigskip
\begin{tabular}{@{}lrlrr@{}}\toprule Sitio & $\varphi$ ($^\circ$) & Pasos cenitales & $N_{jun}$ & $N_{dic}$ \\ \midrule
Humedal La Segua & $-$0{,}70 & 19/03 y 25/09 & 11{,}96 & 12{,}04 \\
Bosques andinos de Loja & $-$4{,}05 & 11/03 y 04/10 & 11{,}77 & 12{,}23 \\
Pedernales & $-$0{,}07 & 21/03 y 24/09 & 12{,}00 & 12{,}00 \\
Glaciares del volcán Antisana & $-$0{,}47 & 20/03 y 25/09 & 11{,}97 & 12{,}03 \\
Arrozales de Daule & $-$1{,}86 & 16/03 y 28/09 & 11{,}89 & 12{,}11 \\
Guayaquil & $-$2{,}19 & 16/03 y 29/09 & 11{,}87 & 12{,}13 \\
\bottomrule\end{tabular}\end{table}}

\newcommand{\TablaVerificacion}{%
\begin{table}[htbp]\centering\footnotesize\setlength{\tabcolsep}{4pt}
\caption{Verificación: indicadores de cada escenario de Ecuador sin forzamiento estacional (fase fija heredada) y con forzamiento climatológico. Semilla 7, 12 \emph{looks}, órbita ascendente.}\label{tab:verificacion}
\begin{tabular}{@{}llll@{}}\toprule Tipo & Sitio & Fase fija & Forzamiento estacional \\ \midrule
Pérdida de humedales & Humedal La Segua & máx. 31/03, mín. 14/11 & máx. 19/03, mín. 26/11 \\
Incendios forestales & Bosques andinos de Loja & amplitud HV 0{,}1 dB & amplitud HV 0{,}1 dB \\
Terremotos & Pedernales & $\sigma_{atm}$ constante, 1{,}80 cm & $\sigma_{atm}$ 1{,}69--1{,}89 cm \\
Movimiento de glaciares & Glaciares del volcán Antisana & retroceso 15 m & retroceso 15 m \\
Actividad agrícola & Arrozales de Daule & exactitud 98,9 \% & exactitud 98,9 \% \\
Subsidencia del terreno & Guayaquil & RMSE 1,98 cm/año & RMSE 1,93 cm/año \\
\bottomrule\end{tabular}\end{table}}

\newcommand{\TablaInvariancia}{%
\begin{table}[htbp]\centering\small
\caption{Prueba de invariancia: fecha del máximo del pulso de inundación en La Segua (sin tendencia de secado) al desplazar la fecha de referencia. Con fase fija el pico se desplaza con la fecha; con forzamiento permanece en la época lluviosa.}\label{tab:invariancia}
\begin{tabular}{@{}lrll@{}}\toprule Fecha de referencia & $\Delta$ (días) & Pico con fase fija & Pico con forzamiento \\ \midrule
01/05/2026 & +0 & 31/03 & 19/03 \\
30/07/2026 & +90 & 29/06 & 25/03 \\
28/10/2026 & +180 & 27/09 & 26/03 \\
\bottomrule\end{tabular}
\par\medskip
\begin{tabular}{@{}lrl@{}}\toprule Fecha de referencia (Daule) & $\Delta$ (días) & Meses de cosecha de arroz \\ \midrule
15/05/2026 & +0 & abr, may, sep, oct \\
12/09/2026 & +120 & abr, may, sep, oct \\
\bottomrule\end{tabular}\end{table}}

\newcommand{\TablaParticion}{%
\begin{table}[htbp]\centering\small
\caption{Resultado de la partición de demostración: 25 bloques de 8~km, buffer de 1500~m y semilla 20261011. Superficie efectiva tras el recorte por buffer.}\label{tab:particion}
\begin{tabular}{@{}lrrrr@{}}\toprule Conjunto & Bloques & Superficie (km$^2$) & \% superficie & Recortado (km$^2$) \\ \midrule
Entrenamiento & 15 & 708{,}91 & 52{,}6 & 251{,}09 \\
Validación & 5 & 320{,}00 & 23{,}7 & 0{,}00 \\
Prueba & 5 & 320{,}00 & 23{,}7 & 0{,}00 \\
Excluido por buffer & 0 & 0{,}00 & 0{,}0 & 0{,}00 \\
\bottomrule\end{tabular}
\par\smallskip{\scriptsize Huella SHA-256: \texttt{c61528b146153d03d755f9ba95832bfd}\allowbreak\texttt{557374fce2bf3af97593705eb1b86d39}}
\end{table}}
\newcommand{\BufferDemo}{1500}
\newcommand{\BloqueDemo}{8000}

\cptitulocorto{Partición espacial reproducible · validación sin fuga}

\begin{document}
\cpportada{> NISAR // VALIDACIÓN // PARTICIÓN ESPACIAL\_}
  {Teoría de la partición espacial reproducible para validar modelos con datos SAR}
  {Bloques, zona de exclusión, autocorrelación y huella criptográfica}
  {Documentación técnica del proyecto \emph{modelo clima}}
  {\today}

\begin{cpcaja}[title=RESUMEN]
Se formaliza el algoritmo de \codigo{modelo/clima.py}, que divide un dominio geográfico en bloques
espaciales y los asigna a entrenamiento, validación y prueba para evaluar clasificadores de cambios
entrenados con píxeles NISAR sin fuga de información espacial. Se describen la discretización del dominio,
la asignación aleatoria reproducible, la zona de exclusión (\emph{buffer}) con recorte de los bloques de
entrenamiento, el criterio para elegir su anchura a partir de la longitud de correlación del campo, su
justificación física mediante la familia de covarianzas de Matérn, las propiedades que verifican las
pruebas automáticas y la huella SHA-256 que certifica la reproducibilidad. Esta versión corrige la
descripción anterior del buffer: los bloques de entrenamiento no se excluyen enteros, sino que pierden
solo la franja situada a menos de $d$ de validación o prueba.
\end{cpcaja}

\tableofcontents
\clearpage

%=====================================================================================================
\section{Introducción}
%=====================================================================================================

Los datos climáticos y de radar no son independientes entre puntos vecinos: dos píxeles cercanos se
parecen mucho más que dos píxeles separados por varios kilómetros. Esta \emph{autocorrelación espacial}
invalida la validación cruzada aleatoria, porque entrenamiento y prueba comparten zonas contiguas y las
métricas resultan optimistas \cite{roberts2017,ploton2020}. La solución estándar es la validación por
bloques espaciales con una separación mínima entre conjuntos \cite{roberts2017,valavi2019}.

El módulo \codigo{modelo/clima.py} implementa esa estrategia con tres objetivos:
\begin{enumerate}
  \item definir una partición espacial reproducible sobre una cuadrícula regular;
  \item garantizar que ningún punto de entrenamiento quede a menos de una distancia $d$ de validación o
        prueba;
  \item certificar la integridad del resultado con pruebas automáticas y una huella criptográfica.
\end{enumerate}
La pestaña \emph{Validación} de la aplicación muestra la partición de demostración
(tabla~\ref{tab:particion}) y sus artefactos se guardan en \codigo{resultados/particion\_espacial/}.

%=====================================================================================================
\section{Dominio espacial y discretización}
%=====================================================================================================

Sea $\Omega\subset\mathbb R^2$ el dominio de estudio, unión de las geometrías del conjunto de datos,
expresado en un sistema de coordenadas \emph{proyectado} con unidades métricas: el código rechaza los
sistemas geográficos (latitud/longitud), en los que la distancia euclídea no tiene unidades de longitud.
Para Ecuador continental suele usarse WGS~84 / UTM~17S (EPSG:32717), con la zona 18 al este de
$78^\circ$\,O y las zonas 15--16 en Galápagos. La demostración usa una malla sintética de
$40\times40$~km con el código EPSG:32630 solo como etiqueta métrica.

Con $(x_{\min}, y_{\min}, x_{\max}, y_{\max})$ la caja envolvente de $\Omega$ y $L$ el lado del bloque,
la cuadrícula es
\begin{equation}
  B_{ij} = [x_i, x_i + L]\times[y_j, y_j + L],\qquad x_i = x_{\min} + iL,\quad y_j = y_{\min} + jL,
\end{equation}
con $i = 0,\dots,\lceil (x_{\max}-x_{\min})/L\rceil - 1$ y análogamente para $j$. Solo se conservan los
bloques que tocan los datos,
\begin{equation}
  \sum_{g}\operatorname{área}(g\cap B_{ij}) > 0,
\end{equation}
donde la suma recorre las geometrías $g$ del conjunto. La superficie que se informa para cada bloque es
la efectiva, $A_{ij} = \operatorname{área}(B_{ij}\cap\Omega)$, reducida por el recorte del buffer en los
bloques de entrenamiento (sección~\ref{sec:buffer}). Se exigen al menos cinco bloques.

%=====================================================================================================
\section{Asignación aleatoria reproducible}
%=====================================================================================================

Con $n$ bloques y la semilla $s$ se inicializa el generador pseudoaleatorio de NumPy y se obtiene una
permutación $\pi = \operatorname{perm}(\{1,\dots,n\};\,s)$. Los tamaños de los conjuntos son
\begin{equation}
  n_{\mathrm{train}} = \lfloor 0{,}6\,n\rfloor,\qquad n_{\mathrm{val}} = \lfloor 0{,}2\,n\rfloor,\qquad
  n_{\mathrm{test}} = n - n_{\mathrm{train}} - n_{\mathrm{val}},
\end{equation}
de modo que el redondeo favorece siempre a la prueba ($n_{\mathrm{test}}\ge 0{,}2\,n$). Los primeros
$n_{\mathrm{train}}$ índices de $\pi$ forman $\trainset$, los $n_{\mathrm{val}}$ siguientes $\valset$ y el
resto $\testset$. Para una misma geometría, $L$, $d$ y $s$, la partición es idéntica en cualquier
ejecución. La proporción 60/20/20 equilibra el ajuste del modelo, la selección de hiperparámetros y la
evaluación final.

%=====================================================================================================
\section{Autocorrelación espacial y anchura del buffer}
%=====================================================================================================

Para un campo $Z(\vect x)$ estacionario, la dependencia entre dos puntos separados por $\vect h$ se
describe con la covarianza $C(\vect h) = \operatorname{Cov}[Z(\vect x), Z(\vect x+\vect h)]$ o con el
semivariograma \cite{matheron1963,cressie1993}
\begin{equation}
  \gamma(\vect h) = \tfrac12\operatorname{E}\!\left[(Z(\vect x+\vect h)-Z(\vect x))^2\right] = C(\vect 0) - C(\vect h).
\end{equation}
En el modelo exponencial, $C(h) = \sigma^2 e^{-h/\ell}$, la correlación entre dos puntos a distancia $h$
es $\rho(h) = e^{-h/\ell}$ y el \emph{alcance práctico}, donde $\gamma$ alcanza el 95~\% de la meseta,
es $a = 3\ell$. Si se exige que la correlación residual entre entrenamiento y evaluación no supere
$\rho_{\max}$, la anchura del buffer debe cumplir
\begin{equation}\label{eq:buffer}
  d \;\ge\; \ell\,\ln\frac{1}{\rho_{\max}},
\end{equation}
y con $\rho_{\max} = e^{-3}\approx 0{,}05$ el criterio coincide con tomar $d\ge a$. Por ejemplo, el
buffer de demostración, $d = \BufferDemo$~m, garantiza $\rho\le 0{,}05$ para campos con
$\ell\le 500$~m. En la práctica, $\ell$ se estima ajustando el semivariograma empírico de los datos (o
de los residuos del modelo) antes de fijar $d$ \cite{roberts2017,valavi2019}.

%=====================================================================================================
\section{Zona de exclusión y recorte del entrenamiento}\label{sec:buffer}
%=====================================================================================================

Sea $V = \bigcup_{B\in\valset\cup\testset} B$ la unión de los bloques de validación y prueba. La zona de
exclusión es el conjunto de puntos a distancia menor que $d$ de $V$:
\begin{equation}
  \mathcal B_d = \{\vect x\in\mathbb R^2 : \operatorname{dist}(\vect x, V) < d\}.
\end{equation}
El código la construye con el \emph{buffer} poligonal de \emph{shapely}, que aproxima cada arco por
$q = 16$ segmentos por cuarto de circunferencia. Un polígono regular de $N = 4q$ lados inscrito en una
circunferencia de radio $R$ tiene apotema $R\cos(\pi/N)$; para que el polígono contenga el disco de
radio $d$ se usa
\begin{equation}
  R = \frac{d}{\cos\!\left(\pi/4q\right)} \approx 1{,}0012\,d,
\end{equation}
de modo que todo punto fuera de la zona construida está a distancia $\ge d$ de $V$.

\paragraph{Recorte, no exclusión.} Cada bloque de entrenamiento conserva solo la parte que queda fuera
de la zona:
\begin{equation}
  E_{k} = (B_k\cap\Omega)\setminus\mathcal B_d,\qquad
  A^{\mathrm{rec}}_{k} = \operatorname{área}(B_k\cap\Omega) - \operatorname{área}(E_k),
\end{equation}
y solo se marca como \emph{excluido por buffer} si $\operatorname{área}(E_k) = 0$. Por construcción,
\begin{equation}
  \operatorname{dist}(E_k,\, V)\ \ge\ d\qquad\forall\, B_k\in\trainset .
\end{equation}
La versión anterior de este documento afirmaba que el bloque se excluía entero en cuanto tocaba la zona.
Ese criterio vacía el entrenamiento: con asignación aleatoria, un bloque interior tiene ocho vecinos (la
zona alcanza también a los diagonales para cualquier $d>0$) y la probabilidad de que ninguno sea de
validación o prueba es $0{,}6^{8}\approx 0{,}017$, de modo que se perdería cerca del 98~\% del
entrenamiento. El recorte conserva la superficie útil sin relajar la garantía de distancia. En la
demostración (tabla~\ref{tab:particion}) ningún bloque queda excluido y el recorte retira una parte de
la superficie de entrenamiento.

\TablaParticion

%=====================================================================================================
\section{Fundamento físico de la correlación espacial}
%=====================================================================================================

Muchas variables climáticas (temperatura, humedad, contenido de agua del suelo) obedecen, a escala
regional, una ecuación de transporte del tipo
\begin{equation}
  \frac{\partial\phi}{\partial t} + \nabla\cdot(\vect u\,\phi) = \nabla\cdot(K\nabla\phi) - \lambda\,\phi + S,
\end{equation}
con advección $\vect u$, difusividad $K$, tasa de pérdida $\lambda$ y forzamiento $S$. En el caso
estacionario, sin advección y con forzamiento aleatorio blanco $W$, el campo satisface
$(\kappa^2 - \nabla^2)\,\phi = W/K$ con $\kappa^2 = \lambda/K$. Whittle \cite{whittle1954} mostró que
las soluciones de ecuaciones de este tipo, $(\kappa^2 - \nabla^2)^{\alpha/2}Z = W$, tienen covarianza de
la familia de Matérn,
\begin{equation}
  C(h) = \frac{\sigma^2}{2^{\nu-1}\Gamma(\nu)}\,(\kappa h)^{\nu} K_{\nu}(\kappa h),\qquad
  \nu = \alpha - \tfrac{d_e}{2},
\end{equation}
donde $K_\nu$ es la función de Bessel modificada de segunda especie y $d_e = 2$ la dimensión del
dominio \cite{lindgren2011}. El caso $\nu = 1/2$ es la covarianza exponencial de la sección anterior, y
el alcance efectivo es $r\approx\sqrt{8\nu}/\kappa$ \cite{lindgren2011}. Así, la longitud de
correlación que fija la anchura del buffer (ecuación~\ref{eq:buffer}) tiene un significado físico: crece
con la difusividad y disminuye con la tasa de pérdida, $\ell\sim\sqrt{K/\lambda}$. La topografía, la
circulación atmosférica y la heterogeneidad de la superficie añaden estructura, pero no eliminan esta
dependencia de corto alcance; por eso la separación entre conjuntos debe ser física y no aleatoria.

%=====================================================================================================
\section{Integridad y pruebas automáticas}
%=====================================================================================================

La partición cumple, por construcción,
\begin{equation}
  \trainset\cap\valset = \trainset\cap\testset = \valset\cap\testset = \varnothing,
\end{equation}
y el módulo incluye cinco pruebas unitarias (\codigo{python -m unittest modelo.clima}):
\begin{enumerate}
  \item \codigo{test\_no\_block\_shared}: ningún bloque aparece en más de un conjunto activo;
  \item \codigo{test\_buffer\_exclusion\_respected}: ningún bloque excluido por buffer queda en entrenamiento;
  \item \codigo{test\_training\_area\_outside\_buffer}: existe entrenamiento y toda geometría efectiva
        $E_k$ está a distancia $\ge d$ de validación y prueba;
  \item \codigo{test\_reproducibility}: dos ejecuciones con la misma semilla producen la misma huella;
  \item \codigo{test\_proportions\_tolerance}: hay al menos un bloque activo.
\end{enumerate}
La última es una comprobación mínima: no contrasta las proporciones obtenidas con las nominales. Estas
pruebas verifican propiedades de la partición, no la similitud estadística entre conjuntos, que es
necesaria pero no suficiente para un buen diseño experimental.

%=====================================================================================================
\section{Huella criptográfica y reproducibilidad}
%=====================================================================================================

La integridad de los artefactos se certifica con una huella SHA-256 \cite{fips180}:
\begin{equation}
  H = \mathrm{SHA256}\!\left(\mathrm{JSON}_{\mathrm{ord}}(C)\;\|\;\mathrm{CSV}\big(\operatorname{ord}_{\mathrm{block\_id}}(\mathcal D)\big)\right),
\end{equation}
donde $C$ es la configuración (sistema de referencia, $L$, $d$, semilla, proporciones y número de
bloques) serializada con las claves ordenadas, $\mathcal D$ la tabla de bloques con sus atributos y
geometrías en WKT, y $\|$ la concatenación. Ordenar las claves y los bloques elimina toda ambigüedad,
de modo que cualquier cambio de semilla, de parámetros o de geometría produce una huella distinta. Los
cuatro artefactos (\codigo{spatial\_split.csv}, \codigo{split\_config.yaml}, \codigo{split\_report.json}
y \codigo{split\_hash.txt}) se escriben en \codigo{resultados/particion\_espacial/}, con independencia
del directorio desde el que se ejecute el módulo.

%=====================================================================================================
\section{Conclusión}
%=====================================================================================================

La partición de \codigo{modelo/clima.py} combina una discretización regular del dominio, una asignación
aleatoria reproducible y una zona de exclusión con recorte que garantiza una distancia mínima $d$ entre
entrenamiento y evaluación sin vaciar el entrenamiento. La anchura $d$ debe elegirse a partir de la
longitud de correlación del campo (ecuación~\ref{eq:buffer}), que tiene una interpretación física a
través de la familia de Matérn. Las pruebas automáticas y la huella SHA-256 hacen la partición
verificable y reproducible.

%=====================================================================================================
\begin{thebibliography}{99}\small

\bibitem{cressie1993} Cressie, N. (1993). \emph{Statistics for Spatial Data}, ed. revisada. Wiley, Nueva
York.

\bibitem{fips180} NIST (2015). \emph{Secure Hash Standard (SHS)}. Federal Information Processing
Standards Publication 180-4. National Institute of Standards and Technology, Gaithersburg.

\bibitem{lindgren2011} Lindgren, F., Rue, H. y Lindström, J. (2011). An explicit link between Gaussian
fields and Gaussian Markov random fields: the stochastic partial differential equation approach.
\emph{Journal of the Royal Statistical Society B}, 73(4), 423--498.

\bibitem{matheron1963} Matheron, G. (1963). Principles of geostatistics. \emph{Economic Geology}, 58(8),
1246--1266.

\bibitem{ploton2020} Ploton, P. et al. (2020). Spatial validation reveals poor predictive performance of
large-scale ecological mapping models. \emph{Nature Communications}, 11, 4540.

\bibitem{roberts2017} Roberts, D.~R. et al. (2017). Cross-validation strategies for data with temporal,
spatial, hierarchical, or phylogenetic structure. \emph{Ecography}, 40(8), 913--929.

\bibitem{valavi2019} Valavi, R., Elith, J., Lahoz-Monfort, J.~J. y Guillera-Arroita, G. (2019). blockCV:
An R package for generating spatially or environmentally separated folds for k-fold cross-validation of
species distribution models. \emph{Methods in Ecology and Evolution}, 10(2), 225--232.

\bibitem{whittle1954} Whittle, P. (1954). On stationary processes in the plane. \emph{Biometrika}, 41(3/4),
434--449.

\end{thebibliography}
\end{document}
